%% =====================================================================
%%  第 5 章　隐式字段
%%  源文件：05-隐式字段.md
%% =====================================================================
\chapter{隐式字段}

\applicable{你收到一条消息或听到一句话，字面上挑不出毛病，但你隐约觉得哪里不对。}

\begin{guideblock}
\textbf{本章解决什么问题}：为什么同样一句“好的”，有时候是同意，有时候是生气；以及怎么读出对方没有说出口的那部分。
\end{guideblock}

先给结论：\textbf{人类接口传输的数据里，只有一部分是明文的。}剩下那部分叫隐式字段——语气、停顿、标点、回复延迟。它们不在字面里，但对方确实发出去了。

\textbf{你把隐式字段读漏了，就等于丢包。}而且这类丢包最麻烦的地方在于：报文本身是完整的，所以你不会察觉自己丢了东西。本章要做的事情只有两件：先把这些没有名字的字段列出来，再给出一套读它们的方法。

\hcirule

\section{同一个“好的”，几种意思}

先看一组样本。

\sample{0927}
\begin{mdquote}
　同一个发送者，对同一件事的四种回复：

　① 好的\par
　② 好的\textasciitilde\par
　③ 好的。\par
　④ 好的！！
\end{mdquote}

字面内容完全一致，都是“我同意了”。

而在绝大多数以中文为母语的人那里，这四句读出来的意思是：

\begin{manstable}{P{2.6cm} L}
\tbh{回复} & \tbh{通常的含义} \\
\midrule
\textbf{好的} & 中性。同意，没有更多信息 \\
\textbf{好的\textasciitilde} & 偏暖。同意，且愿意接这个事 \\
\textbf{好的。} & 偏冷。同意，但\textbf{加了一个句号}，说明态度被刻意收了回去 \\
\textbf{好的！！} & 异常。要么真的很高兴，要么在生气——\textbf{方向取决于这个人平时发不发感叹号} \\
\bottomrule
\end{manstable}

需要强调的是，上面这张表\textbf{不是一个可靠的字典}。它只在一种条件下成立：你和这个人已经建立了基线。而基线的建立方式，是 5.4 的内容。

值得注意的是第④种。\textbf{它不是“更热情”，它是“偏离”。}而偏离本身才是信号——至于往哪个方向偏离，要靠 5.4 的方法判断。

（同一个词，来自不同的人，意思不一样。先看这个对照。）

上面那张表说的都是“某一个固定的人”。把它换个人，结论会变。看下面这个样本。

\sample{0950}
\begin{mdquote}
　同一天，同一个词“好的”，来自三个不同的人：

\textbf{发消息的人一：你的室友。}平时回你消息从来不打标点，也不打“好”，都是“行”“OK”“来了来了”，平均回复时间约 4 分钟。

　今天他回的是：「好的。」

\textbf{发消息的人二：你妈。}平时每句话后面都带感叹号或者表情包，平均回复时间约 20 分钟。

　今天她回的是：「好的。」

\textbf{发消息的人三：你合作了三年的同事。}平时说话一板一眼，每句话都带句号，连“收到。”都要加句号。

　今天他回的是：「好的。」
\end{mdquote}

同一个词，三个发送者，三种含义。

\begin{itemize}
\item \textbf{室友那个句号，是加出来的。}他从不用标点，这次的句号是他特意打上去的。\textbf{加了一个原本不存在的东西，等于在报文体里插了一行注释。}含义通常是：他不太想做这件事，或者他有点不高兴，但不想说。
\item \textbf{你妈那个句号，是减出来的。}她平时不用句号，是因为她习惯用感叹号表达热心。今天换成句号，说明热心的那部分被收回去了。含义通常是：有点失落，或者有点不高兴。
\item \textbf{同事那个句号，什么都不是。}那是他的基线，他平时就这样。\textbf{在他身上读句号，等于把正常协议头当成错误。}
\end{itemize}

\textbf{这就是为什么本章不给“标点对照表”。}同一串字节，在不同的人身上解出来是不同的东西。你只能先知道这个人平时什么样。

还有一层：\textbf{这三个人的回复时间也都在传信息。}室友平时 4 分钟回你，这次如果拖到 40 分钟，那是异常；你妈平时 20 分钟，这次 4 分钟回了个句号，那也是异常。偏离可以是标点上的，也可以是时间上的，两者常常同时出现。\textbf{同时出现两次偏离的时候，信号强度远高于一次。}

\section{隐式字段清单}

\hciTag{核心结论}：人类一句话里，除字面外，还有这些字段同时在传。它们都没有名字，所以只能靠归纳。

\begin{manstable}{P{2.2cm} P{1.6cm} L}
\tbh{字段} & \tbh{载体} & \tbh{读什么} \\
\midrule
\textbf{语气} & 声音 & 是升调还是降调；是软的还是硬的 \\
\textbf{语速} & 声音 & 变快通常是紧张或敷衍；变慢通常是沉重或试探 \\
\textbf{停顿} & 声音 & \textbf{停在哪里，比说了什么更重要} \\
\textbf{重音} & 声音 & 加重的那两个字，是真正想说的 \\
\textbf{音量} & 声音 & 和平时比是偏高还是偏低；偏低有时候比偏高更值得注意 \\
\textbf{叹气} & 声音 & 出现在句首还是句尾，含义不同 \\
\textbf{面部表情} & 面部 & 与内容是否一致（说“没事”却在抿嘴） \\
\textbf{视线} & 面部 & 说话时看着你还是看着别处 \\
\textbf{身体朝向} & 姿态 & 朝向你是开放，侧过去是收缩 \\
\textbf{手部动作} & 姿态 & 双臂交叉、翻手机、反复看表 \\
\textbf{标点} & 文字 & 句号、波浪号、感叹号的增减 \\
\textbf{回复延迟} & 时间 & \hciSee{5.3} \\
\textbf{回复长度} & 文字 & 平时三行的人回了一个字；平时一个字的人回了三行 \\
\textbf{渠道选择} & 媒介 & 打电话 vs 发微信 vs 当面，本身就是一次表态 \\
\textbf{发送时段} & 时间 & 深夜的消息和下午三点的消息，重量不同 \\
\bottomrule
\end{manstable}

其中语气、语速、停顿、重音这几项，你在电话里能全部拿到，在文字里全部拿不到。\textbf{这就是文字消息比通话容易误读的物理原因。}

但在文字里，你也不是完全没有东西可读。文字里可读的部分是——\textbf{标点、回复延迟、回复长度、渠道、时段}。这五项是文字通道里仅剩的隐式字段。\textbf{把消息发成纯文字的人，等于把这五项以外的东西全部丢掉了。}

\section{回复延迟是最被忽略的字段}

\hciTag{注意事项}：这是本章最实用的一节。

在人类接口里，\textbf{“多久回复”本身就是一个字段，而且它是自动发送的}——你无法选择不发送它。哪怕你什么都不回，那个沉默的时长也在传递信息。

一条消息发出后，回复延迟通常被这样解读：

\begin{manstable}{P{2.6cm} L}
\tbh{延迟} & \tbh{通常的解读} \\
\midrule
秒回 & 重视，或者正好在玩手机 \\
几十分钟 & 正常 \\
几小时 & 可能在忙，也可能在犹豫 \\
隔天 & 信号很强 \\
不回 & \hciSee{第6章}的“没事”——沉默是一种返回码 \\
\bottomrule
\end{manstable}

需要注意两点。

\textbf{第一，延迟对照的是基线，不是绝对值。}一个平时三天不回消息的人，两小时回你，那是“非常快”。一个平时秒回的人，隔了两小时，那是“出事了”。\textbf{看绝对值毫无意义。}

\textbf{第二，延迟读的是变化，不是长度。}同一个人回复你的速度突然变慢，比他一直很慢，信息量大得多。

关于延迟，还有三点补充，都是实战中真正会遇到的情形。

\textbf{第三，延迟的长度和事情的重要性不成正比。}很多人以为“重要的事应该秒回”，实际观测相反：\textbf{一条越需要斟酌的回复，延迟越长。}回复“今天去哪吃”可以秒回，回复“你觉得我们这样下去行吗”往往要隔几小时。\textbf{所以延迟变长，有时候不是冷淡，而是在认真组织内容。}这两种情况无法从延迟本身区分，只能靠后续那条消息的质量来判断。

\textbf{第四，延迟的计算基准是“这条链路”，不是“这一天”。}一个平时秒回、今天下午隔了三小时才回的人，那三小时是信号。但要注意，如果这三小时里他也在给别人秒回，那不成立——你看不到这一层，所以不要往这个方向使劲。

\textbf{第五，不要统计短窗口。}只看最近三次消息的延迟，得出的结论基本是噪声。\textbf{至少要看两周，才谈得上基线。}

\section{读隐式字段的方法：先基线，再偏离}

这是本章唯一的方法论，也是全书反复要用的一个思维工具。

\textbf{第一步，建立基线。}观察这个人平时怎么说话：他平时用不用句号？他平时多久回一条消息？他平时打电话是什么语速？\textbf{基线不是“一般人怎么样”，是“他这个人怎么样”。}

基线要记的东西不多，只有五条：

\begin{manstable}{P{2.6cm} L}
\tbh{基线条目} & \tbh{记什么} \\
\midrule
标点习惯 & 用不用句号、用不用波浪号、用不用表情 \\
回复节奏 & 通常多久回；白天回得快还是晚上回得快 \\
句子长度 & 通常是一次说一句，还是连着发好几条 \\
语音语速 & 打电话时快不快、停顿多不多 \\
开场收尾 & 会不会打招呼，会不会说“先这样” \\
\bottomrule
\end{manstable}

\textbf{第二步，看偏离。}找出与基线不一致的地方。

\begin{mdquote}
一个从不发句号的人，突然发了一个句号。\par
一个秒回的人，隔了一天才回。\par
一个说话很慢的人，突然语速变快。\par
一个每次都要寒暄两句的人，这次直接进正题。
\end{mdquote}

\textbf{偏离就是信号。}而且偏离不需要你判断方向——你只需要知道“这里有东西”，然后去 5.5 确认。

\textbf{偏离分三档，反应方式完全不同。}

\begin{manstable}{P{2.2cm} L P{3.4cm}}
\tbh{强度} & \tbh{表现} & \tbh{该怎么办} \\
\midrule
\textbf{单次偏离} & 只在一个字段上不同于平时 & 记下，不反应。可能只是手滑 \\
\textbf{连续偏离} & 两个以上字段同时不同，或同类偏离出现三次以上 & 当疑问处理，找一个低压力的时刻问 \\
\textbf{持续偏离} & 超过一周都不同，或同时伴随渠道改变（从语音变文字） & 当状态处理，不再是这一条消息的事 \\
\bottomrule
\end{manstable}

\textbf{第三，不要用偏离去定罪。}偏离只说明“发生了异常”，不说明“异常是什么”。一个变慢的回复，可能是生气，也可能是他手机坏了。\textbf{把偏离当成疑问，不要当成结论。}

\textbf{一个容易忽略的点：偏离也可能来自好事。}一个人突然开始主动发消息、突然开始用表情、突然回复变快——这些同样是偏离，方向是正的。\textbf{好消息和坏消息都以偏离的形式到达，而好消息更容易被当成理所当然忽略掉。}

\section{共情：读对方没有发送的字段}

第 9 章会详细讲“接情绪”。这里先给一个定义，因为它属于本章的范围。

\textbf{共情，是解析对方没有发送的字段。}

它不是“我也有过同样的经历”——那是把自己的数据包发回去。它是：对方只说了一句话，而你把这句话背后那部分也读进来了。

\begin{mdquote}
对方说：“没事，你忙你的。”

\textbf{只读字面}：他让我去忙，那就去忙。\par
\textbf{读完隐式字段}：语气是平的，标点是句号（他平时不用句号），回复比平时快——\textbf{偏离了三次。}\par
\textbf{共情}：他知道我有事，不想占用我，但他自己那边有事。
\end{mdquote}

共情的输出不是分析，是一句话：

\begin{mdquote}
“我这个事不急。你那边怎么了？”
\end{mdquote}

\textbf{注意这句话的结构：它没有追问（“你到底怎么了”），也没有否认对方的说法（“你明明有事”），它只是把对方的“你忙你的”往后放了一位。}

关于共情还有一条边界，需要说清楚。

\textbf{读得出来，不等于说得出来。}有的人察觉到了对方的偏离，第一反应是把它讲出来：“你今天语气不对”“你今天怎么不用波浪号了”。\textbf{这句话的问题在于，它把对方隐藏的字段摊到了桌面上。}对方的反应通常是收回、防御，或者否认。\textbf{察觉是接收，说出来是发送。}这两件事要分开。

\section{读错隐式字段的三种典型代价}

这一节讲读错的后果。之所以要单独讲，是因为隐式字段的误读有一个特征：\textbf{它当场不报错。}

你读错了，对方不会说“你读错了”。对话照常进行，事情照常往下走。代价要过一段时间才出现，而且出现的方式往往和你读错的那件事没什么表面关系。

\textbf{代价一：把沉默读成否定}

这是最高频的一种，也是最容易在文字通道里发生的。

\sample{0953}
\begin{mdquote}


周三晚上九点，你给一个合作方发了一条微信：「张老师，方案我这边周五能出，您看节奏合适吗？」

消息发出去，显示已读，没有回复。

晚上十一点，你又看了一眼，还是没有回复。

你开始重写了一版方案，把不太确定的部分全删了，改成保守版本。凌晨一点你又看了两次手机，仍然没有回复。

第二天早上十点，对方回了一句：「合适，麻烦你了。」
\end{mdquote}

这条消息从头到尾没有任何恶意。对方只是晚上没看工作消息。

但你在这 13 个小时里做的事情是：\textbf{把一段没有内容的沉默，当成了一个否定的回答。}你重做了方案，压低了风险，多花了一个晚上。\textbf{沉默是空包体，它不携带内容。}把它填上你自己最担心的那个答案，是人在等消息时的默认行为。

代价不是这一次的加班。\textbf{代价是下次你还会这么做}——而且你会因为“上次也是这样”而更加确信自己的解读是对的。

\textbf{代价二：把偏离读成都市传说}

这一种发生在你想得太多的时候。

\sample{0954}
\begin{mdquote}
　你和一个人连续三周每天聊几句。这三周里，她回复你一般在一小时内，会带表情，会说“嗯嗯”。

今天她回你：「嗯。」

你注意到这个“嗯”没有重复、没有表情、多等了 40 分钟。

你开始回想：是不是上周五那句话说得不对？是不是她最近在忙？是不是她不想理我了？

你翻了最近一周的聊天记录，找出了三个“可疑的点”。你决定这两天先不主动找她，看看她会不会先开口。
\end{mdquote}

三天后，你从共同朋友那里听说：她那几天在赶一个项目的截止日期，跟谁说话都是“嗯”。

\textbf{这个样本里，你真正读错的地方不是那一个“嗯”。}你读对了一半——“嗯”确实是偏离。\textbf{你读错的是把一次偏离，升级成了一条推论链。}偏离只说明“这里有东西”，它不说明是什么。一次偏离的正确处理方式是记下，不是启动调查。

值得注意的是，\textbf{这种代价比第一种更隐蔽}，因为它表现为“你在乎”。人在过度解读的时候，自我感觉是“我很用心”。而在对方那边，体验是“这个人怎么突然不理我了”。

\textbf{代价三：把基线读成通用的}

这一种代价最大，因为它会稳定地、反复地污染你所有的判断。

\sample{0955}
\begin{mdquote}
　你有两个朋友。

\textbf{朋友甲}：回消息很快，平均十分钟，句尾爱用波浪号，每次都带表情。\par
\textbf{朋友乙}：回消息很慢，平均一天，从不带表情，从来都是干巴巴一句话。

某天你用你从朋友甲身上学来的标准，去看朋友乙的一条回复：「知道了。」

按朋友甲的标准，“知道了”三个字加句号，是明显的冷。按朋友乙的标准，这是他今天说过的最长的一句话。
\end{mdquote}

\textbf{这类错误的根源，是把自己的习惯当成世界的默认值。}一个秒回的人，容易把所有人的慢都读成不重视；一个说话含蓄的人，容易把所有人的直接都读成冒犯。\textbf{每一次这样的读取，你都在用别人的基线去解别人的报文。}

需要说明的是，这三类代价不会同时发生在一个人身上。多数人只会稳定地犯其中一类。\textbf{先确定你常犯的是哪一类，比同时防三类更有效。}

\section{语音和文字：两套通道的误差对照}

前面一直在说“文字容易丢字段”。这里把两种通道摆在一起，具体对比一次。

先看一个能直接对照的样本。

\sample{0951}
\begin{mdquote}
　同一件事，同一个人，用两种方式说给你听：

\textbf{版本 A（语音）}，时长约 6 秒，语气平，语速正常：「这样吧，那就先这样。」\par
\textbf{版本 B（文字）}：「这样吧，那就先这样。」
\end{mdquote}

两个版本的字面一模一样。但它们携带的字段数量差别很大。

\begin{manstable}{L C C}
\tbh{} & \tbh{语音版本} & \tbh{文字版本} \\
\midrule
语气（升调/降调） & \textbf{有} & 无 \\
语速 & \textbf{有}（6 秒说 9 个字，偏慢） & 无 \\
停顿位置 & \textbf{有}（“这样吧”后面有一个明显的顿） & 无 \\
重音 & \textbf{有}（重音落在“先这样”上） & 无 \\
标点 & 无 & \textbf{有}（“这样吧”后面的逗号是刻意加的） \\
延迟 & \textbf{有} & \textbf{有} \\
回复长度 & 不适用 & \textbf{有}（9 个字，短于他平时） \\
\bottomrule
\end{manstable}

\textbf{语音版本里，“这样吧”后面那个停顿，是整句话唯一真正携带信息的地方。}它说明对方在说“这样吧”之前犹豫了一下。这个停顿在文字版本里完全不存在——文字版本里对应它的只有一个逗号，而逗号是所有标点里信息量最低的一个。

还有一项差别，比字段数量更重要：\textbf{语音是同步的，文字是异步的。}

同步意味着你没有时间准备。对方说完，你的反应立刻发生，\textbf{反应本身就是数据}。异步意味着你有时间组织——文字消息是编辑过的版本。\textbf{你在文字里看到的，是对方想让你看到的那个版本；你在语音里看到的，是对方的实时输出。}

由此得到一个可以直接使用的结论：

\textbf{发生冲突、发生误会、涉及情绪的对话，尽量走语音或当面。}不是因为语音更礼貌，而是因为语音里可读的字段多一倍。\textbf{文字通道里可读的隐式字段，只剩标点、延迟、长度、渠道、时段这五项。}

\section{当对方是上级：压缩最厉害的一种}

职场里的隐式字段有一个特点：\textbf{它被压缩得最厉害，但权重最高。}原因不难理解——在工作关系里，双方都不能把话说满，于是所有多余的信息都被挤进了那些剩下的字段里。

所以读职场信号的方法，和读朋友的不一样。\textbf{对朋友，你读“他多说了什么”；对上级，你读“他比平时少说了什么”。}

\sample{0952}
\begin{mdquote}
　同一个上级，对同一份汇报的三种回复。他平时的基线是：会先回一个“收到”，后面跟一到两句具体意见，通常在 20 分钟内回。

\textbf{回复 A}：「收到，这个方向可以，第三页的数据再核一下。」（18 分钟后）\par
\textbf{回复 B}：「收到。」（18 分钟后）\par
\textbf{回复 C}：「嗯。」（第二天上午 10 点）
\end{mdquote}

三条都是正向或者中性的回复，没有一条带否定词。但按偏离来读：

\begin{manstable}{P{1.4cm} L P{3.6cm}}
\tbh{回复} & \tbh{相对基线的偏离} & \tbh{通常的含义} \\
\midrule
\textbf{A} & 无偏离 & 正常。事在往前走 \\
\textbf{B} & 少了具体意见 & 事放在了“待定”区。\textbf{可能有意见，但不想现在说} \\
\textbf{C} & 晚了约 14 小时，且长度从一句压到一个字 & 态度保留。这个“嗯”不是同意，是“我看到了，先这样” \\
\bottomrule
\end{manstable}

\textbf{关键在于第三条。}它没有否定词，没有感叹号，没有句号——\textbf{它唯一携带信息的地方是：它比他平时短了太多，也晚了很多。}一个平时给具体意见的人，这次只回一个字，说明这次他不想给意见。至于为什么不想给，可能是他还在等别的信息，也可能是他不认可但不想在会上说。

\textbf{这里有一个必须提醒的地方：这个“嗯”不要拿去解读，要拿去做事。}读职场隐式字段的正确用法，是把“他比平时少说了什么”当成一个待办项——\textbf{在你下一次汇报时，用结果把这个缺口补上。}而不是去找他问“您是不是对我有意见”。后者会让他把那个“嗯”变成一个明确的答复，而那个答复大概率比“嗯”更糟。

职场里权重最高的三个位置，按顺序是：

\begin{manstable}{P{2.2cm} L L}
\tbh{位置} & \tbh{基线对照点} & \tbh{信号强度} \\
\midrule
\textbf{回复延迟} & 他通常多久回你 & 高。隔了很久的一个字，通常意味着态度保留 \\
\textbf{回复长度} & 他通常回你多少字 & 高。平时解释两句，这次只回四个字 \\
\textbf{渠道} & 他通常用什么方式找你 & 中。平时发消息，这次直接打电话，事情级别变了 \\
\bottomrule
\end{manstable}

\textbf{“他比平时少说了什么”，是职场里最可靠的信号。}因为在那套环境里，多说一个字都是有成本的，所以少说一个字也是。

\section{什么时候不要读隐式字段}

前面八节都在讲怎么读。这一节讲什么时候不该读。

\textbf{第一条：对陌生人，只读字面。}你和对方没有基线，任何解读都是投射。同事在电梯里没接你的话，可能是他在想事，也可能只是没听见。\textbf{没有基线的解读，读出来的都是你自己。}

\textbf{第二条：一次偏离不构成信号。}这一条在前面出现过，这里再强调一次，因为它是最容易被违反的一条。\textbf{单次偏离的正确反应是“记下”，不是“处理”。}

\textbf{第三条：不要读对方不愿意给的字段。}有些字段是对方刻意收起来的。你读出来了，对方也有权不承认。\textbf{读出来是能力，说出来是另一种行为}，这两者的边界要自己守住。

\textbf{第四条：你自己的状态是最大的噪声源。}你刚吵完架、刚被批评、刚没睡好——在这些状态下，你读到的“偏离”会明显增多。\textbf{这不是对方的信号变多了，是你的接收机灵敏度被调高了。}判断自己当前状态的一个简单办法：过去两小时里，你觉得有几个人对你态度不对？\textbf{超过两个，先怀疑自己。}

\section[常见误区]{常见误区}

\hcimistake{只看字面}{\hciTag{反例}「他说他不生气啊。」}{字面是明文，隐式字段才是大部分载荷。}

\hcimistake{过度解读}{\hciTag{反例}「他一小时没回我，是不是生我气了。」}{一次偏离只是异常，不是结论。慢 20 分钟可能只是下楼拿了快递。}

\hcimistake{把“忙”当“不想理”}{\hciTag{反例}「他以前都秒回的，现在两小时才回，是不想理我了吧。」}{延迟要对照他这两周的基线，不是对照他今天的状态。}

\hcimistake{用自己的基线}{\hciTag{反例}「我都是秒回的，他隔一天才回，这也太不尊重人了。」}{你的习惯不是世界的默认值。他平时就是隔天回。}

\hcimistake{只读单条}{\hciTag{反例}「他就说了句‘好的’，肯定是敷衍我。」}{这一句要放进上下文里。前一小时你们还在正常聊天，那就不是敷衍。}

\hcimistake{把共情做成分析}{\hciTag{反例}「你其实很在意这件事，对吧？」}{这是审讯，不是共情。共情只给一句，然后停下。}

\hcimistake{发现偏离就摊牌}{\hciTag{反例}「你今天怎么怪怪的？」「怎么突然这么客气？」}{会把对方推得更远。多数换来“没有啊”。}

\hcimistake{把字段讲出来}{\hciTag{反例}「你怎么发了个句号？平时你不发的。」}{察觉是接收，说出来是发送。等于把对方的隐藏字段摊在桌上。}

%% ---------- 本章小结 ----------
\hciblocktitle{bullseye}{本章小结}

综上所述，本章的核心结论可以概括为以下几点：

\begin{enumerate}[label=\bfseries\arabic*., leftmargin=2.4em, labelsep=0.7em,
                  itemsep=0.45em, topsep=0.2em, parsep=0.2em]
\item 人类接口的载荷里，只有一部分是明文的，其余是隐式字段。
\item 语气、停顿、标点、回复延迟，都在传东西，而且你无法选择不传。
\item 读隐式字段的方法只有两步：\textbf{先建立这个人的基线，再看偏离}。基线有五条：标点、节奏、长度、语速、开场。偏离有三档：单次、连续、持续。
\item 偏离说明“有异常”，不说明“异常是什么”。把它当疑问，不要当结论。
\item 文字通道里可读的字段只剩五项：标点、延迟、长度、渠道、时段。\textbf{涉及情绪的事，尽量走语音。}
\item 对上级要读“他比平时少说了什么”，而且读出来的东西要拿去做事，不要拿去解读。
\item 共情是解析对方没有发送的字段，输出是一句话，不是一次分析。
\end{enumerate}

%% ---------- 练习 ----------
\begin{exercisessec}
\item 挑一个你每天都会联系的人，写下他的基线：他平时用不用句号、大概多久回消息、一次说一句还是连着发好几条、打电话时语速如何。（[请在此处填写这个人的称呼和你的观察]）

\item 在接下来的一周里，记录三次“对方偏离了基线”的时刻，以及你当时的解读。\textbf{先写下解读，再去验证它对不对。}

\item 把一段最近的语音对话转写下来。对比一下：语音里那些语气、停顿、重音，转成文字之后还剩多少。\textbf{这一步需要你找一个真人，让他对着你说一段 30 秒的话，你先听一遍，再只看转写文字读一遍。}

\item 把你最近发出的一条消息调出来，检查它携带的隐式字段：标点是什么、发出去的时间是几点、延迟了多久才发。判断你在那条消息里，无意中发出了什么。
\end{exercisessec}

%% ---------- 自测题 ----------
\begin{quizsec}
\item 下面哪一项不属于隐式字段？

\begin{enumerate}[label=\Alph*., leftmargin=2.2em, labelsep=0.6em,
                  itemsep=0.2em, topsep=0.3em, parsep=0.1em]
\item 回复延迟
\item 语气
\item 句子里的名词
\item 标点符号
\end{enumerate}

\item 判断：一个人回消息慢，说明他不重视你。

\item 简答：读隐式字段时，“基线”指的是什么？
\end{quizsec}

%% ---------- 参考答案 ----------
\begin{answerssec}
\item C。句子里的名词属于明文内容；延迟、语气、标点都是自动发送的隐式字段。\hciSee{5.2} 与 5.3。

\item 错。延迟要对照这个人自己的基线，绝对速度没有意义。一个平时三天不回消息的人，两小时回你，那是非常快。\hciSee{5.3}。

\item 基线指的是“这个人平时怎么样”，不是“一般人怎么样”。它需要单独建立——他平时用不用句号、多久回一条消息、说话什么语速。没有基线，任何对隐式字段的解读都是投射。\hciSee{5.4}。
\end{answerssec}

\hcirule

\begin{qabox}
\qaQ{读者提问}{如果对方就是不发任何隐式字段，只发干巴巴的事实，怎么办？}
\qaA{本手册回答}{这是一个非常好的问题。需要说明的是，不发送隐式字段本身也是一个字段——它通常说明对方要么在刻意保持距离，要么把这次对话定义为任务型。此时不要追加解读，回到字面，按任务处理。}
\end{qabox}

\hcirule

希望本章的内容能对你有所帮助。如需进一步了解第 6 章《“没事”》的相关内容，我可以随时为你展开。

%% 章末「页脚交互条」由 hci-manual.cls 自动挂接，无需手写。
